% Part: computability
% Chapter: computability-theory
% Section: def-functions-self-reference

\documentclass[../../../include/open-logic-section]{subfiles}

\begin{document}

\olfileid[pt]{cmp}{thy}{slf}
\olsection{Definição de funções por autorreferência}

Em geral, é útil poder definir funções em termos delas próprias. Por
exemplo, dadas funções computáveis $k$, $l$ e~$m$, o lema do ponto
fixo nos diz que existe uma função parcialmente computável~$f$ que
satisfaz a seguinte equação para todo~$y$:
\[
f(y) \simeq
\begin{cases}
k(y) & \text{se $l(y) = 0$} \\
f(m(y)) & \text{caso contrário.}
\end{cases}
\]
Mais especificamente, $f$~é novamente obtida tomando
\[
g(x,y) \simeq
\begin{cases}
k(y) & \text{se $l(y) = 0$} \\
\cfind{x}(m(y)) & \text{caso contrário}
\end{cases}
\]
e usando então o lema do ponto fixo para encontrar um índice~$e$ tal
que $\cfind{e}(y) = g(e,y)$.

Como exemplo concreto, a função ``máximo divisor comum''
$\fn{gcd}(u,v)$ pode ser definida por
\[
\fn{gcd}(u,v) \simeq
\begin{cases}
v & \text{se $u = 0$} \\
\fn{gcd}(\fn{mod}(v, u), u) & \text{caso contrário}
\end{cases}
\]
em que $\fn{mod}(v, u)$ denota o resto da divisão de $v$
por~$u$. Um apelo ao lema do ponto fixo mostra que $\fn{gcd}$ é
parcialmente computável. (Na verdade, isso pode ser colocado no
formato acima, fazendo $y$ codificar o par $\tuple{u, v}$.) Uma indução
subsequente sobre~$u$ mostra, então, que $\fn{gcd}$ é, de fato, total.

Naturalmente, é possível criar definições autorreferenciais muito mais
sofisticadas do que os exemplos que acabamos de discutir. A maioria
das linguagens de programação permite, de uma maneira ou de outra,
definir funções em termos delas próprias. Observe que isso é um pouco
menos impressionante do que poder definir uma função em termos de um
\emph{índice} para um algoritmo que calcula a função, que é o que o
teorema do ponto fixo permite fazer em toda a sua generalidade.

\end{document}
